%%%% Appendices A-F for the TAE compressed version.
%%%% Governing principle: everything cut from the body lives here in full. A reviewer
%%%% reading only the body should be convinced; a reviewer reading the appendices should
%%%% find every body claim substantiated with complete evidence.

\section{Phase 1 in full}
\label{app:phase1}

\paragraph{Aggregation conventions.} A configuration's return is the IQM across seeds of each
seed's mean eval return, with a 95\% stratified-bootstrap CI. The bootstrap is seeded, so
every interval reported in this paper regenerates exactly from the released per-seed data.
Per-room success-rate aggregates are the mean across seeds of each seed's own SR fraction, a
seed-mean rather than a pooled per-episode rate: because eval target rooms are sampled rather
than balanced, seeds carry unequal per-room episode counts, and the seed-mean keeps every
aggregate arithmetically consistent with the per-seed tables (averaging a per-seed column
reproduces its aggregate).

\paragraph{Evidence tiers.} We distinguish three grades of evidence and never let them mix.
Multi-seed Phase 1 and multi-seed Phase 2 results are claim-grade: ten seeds, IQM with 95\%
stratified-bootstrap CIs, and released per-seed data. Single-seed gates are design and
calibration steps only, deciding which rungs are worth spending seeds on, and are never cited
as results or entered into any table.

\paragraph{The nine configurations, and why DQN is among them.} The two architectures compared
in the body are endpoints of a four-point PPO sweep over the actor:critic capacity ratio at
fixed budget. The other five configurations are DQN and carry no asymmetry signal at all. They
are present because Phase 1 holds \emph{two} claims from the earlier work to account, and the
second is a PPO-vs-DQN sample-efficiency comparison, so the five DQN runs are the value-based
comparator family for that claim, not part of the asymmetry test. The PPO configurations vary
essentially only the actor:critic ratio; the DQN configurations (a single Q-network, no
actor--critic split) vary width, discount, exploration schedule, and buffer, and include two
deliberately destabilised settings retained to replicate the original DQN table faithfully.
All nine share the frozen environment, the 200k-step budget, and 10 seeds.

\begin{table}[h]
\caption{The nine Phase-1 configurations. PPO spans the actor:critic ratio; DQN is the
value-based comparator family for the sample-efficiency claim.}
\label{tab:configs}
\centering
\small
\begin{tabular}{llll}
\toprule
Family & Config & Network (actor $\pi$ / critic $v$; DQN: Q-net) & Ratio / role \\
\midrule
PPO & Exp 1 & $\pi$ [256,256] / $v$ [256,256] & symmetric 1:1, the control \\
 & Exp 2 & $\pi$ [256,256] / $v$ [512,256] & conventional (wider critic) \\
 & Exp 3 & $\pi$ [128,128] / $v$ [512,512,256] & conventional (deep critic) \\
 & Exp 4 & $\pi$ [512,256,128] / $v$ [256,128] & \textbf{inverted} (deeper actor) \\
\midrule
DQN & Exp 1 & Q [256,256] & slow $\epsilon$-decay \\
 & Exp 2 & Q [512,256], $\gamma=0.995$ & discount variant \\
 & Exp 3 & Q [512,256] & aggressive $\epsilon$-decay, high lr \\
 & Exp 4 & Q [1024,512] & low lr, high $\gamma$ (unstable) \\
 & Exp 5 & Q [512,256], $\gamma=0.99$ & original best DQN \\
\bottomrule
\end{tabular}
\end{table}

\paragraph{Why the episodes-to-convergence ratio was the wrong measurement.} The earlier work
reported PPO as ${\approx}12.5\times$ more sample-efficient than DQN, on the basis that PPO
reached stable high performance in roughly 20 episodes against nearly 250 for DQN, i.e.\
$12.5\times = 250/20$, a ratio of episodes-to-visual-convergence read off single-run learning
curves. This is not a sound cross-algorithm measure, for three reasons. \textbf{(a) Episodes
are not a common currency of experience.} Episode length on this task varies five- to
ten-fold with policy quality: a wandering early-training agent runs to the 150-step timeout
while a converged agent finishes in ${\sim}14$ steps, so one episode purchases very different
amounts of environment interaction at different points in training. For two algorithms whose
episode-length trajectories differ, an episode-count ratio conflates sample efficiency with
episode duration. The currency the agent actually spends, and that both algorithms share, is
environment steps. \textbf{(b) ``Stable high performance'' is eyeballed}, with no defined
threshold and, read off a single curve, no uncertainty. \textbf{(c) It is a single run.} Our
own H1 data shows how misleading single runs are here: one inverted-configuration seed
(Exp~4, seed~8) collapses to a return of 3.6, a collision local optimum in which terminating
the episode early is easier to find than reaching a distant goal, while its nine siblings sit
at ${\sim}27.6$ and the symmetric configuration has no collapsed seed at all. That asymmetric
collapse is the only asymmetry-shaped signal anywhere in the H1 data.

\paragraph{Re-measured, step-based, multi-seed.} Using environment steps to reach 90\% of the
asymptotic eval-return IQM, over 10 seeds, no advantage remains. To keep the comparison from
turning on a chosen exemplar we report the full within-family distribution rather than a
single pair. The four PPO configurations reach threshold at IQMs of 36.7k (Exp 1), 36.7k
(Exp 2), 163.3k (Exp 3), and 32.0k (Exp 4) steps; the five DQN configurations at 21.7k, 20.0k,
20.0k, 45.0k, and 26.7k, with overlapping confidence intervals throughout (e.g.\ PPO Exp 1
[26.7k, 40.0k], DQN Exp 2 [14.0k, 74.0k]). The two families occupy the same
${\approx}20$--$45$k band, and the single configuration that leaves it, PPO Exp 3 at 163.3k
steps, is the \emph{slowest} of either family. On an exemplar-independent reading PPO is no
faster than DQN, and the one salient difference runs opposite to the claimed $12.5\times$.

\begin{figure}[h]
\centering
\includegraphics[width=0.66\linewidth]{p1_iqm_main_configs.png}
\caption{Per-configuration IQM eval return with 95\% stratified-bootstrap CIs across the nine
main configurations, 10 seeds each.}
\label{fig:p1iqm}
\end{figure}

\begin{figure}[h]
\centering
\includegraphics[width=0.55\linewidth]{p1_perf_profile_ppo_inverted_vs_symmetric.png}
\caption{Performance profile: fraction of runs exceeding a return threshold, inverted (Exp 4)
versus symmetric (Exp 1). The two curves overlap across the range.}
\label{fig:p1prof}
\end{figure}

\begin{figure}[h]
\centering
\includegraphics[width=0.66\linewidth]{p1_sample_efficiency_ppo_vs_dqn.png}
\caption{Sample efficiency: environment steps to 90\% of asymptotic IQM (IQM $+$ 95\% CI),
PPO versus DQN. The two families' distributions overlap; the one outlier is a \emph{slow} PPO
configuration.}
\label{fig:p1sampeff}
\end{figure}

\clearpage
\section{Phase 2 in full}
\label{app:phase2}

\paragraph{Why flicker rate is not the difficulty knob (calibration).} A ground-truth BFS on
the map gives optimal path lengths $H = \{\text{kitchen }11, \text{bedroom }11,
\text{bathroom }22\}$, and an exact consecutive-mask survival model predicts blackout
probability rising with $H$. The observed per-room difficulty contradicts this: under flicker
the two equidistant rooms diverge sharply (at $p=0.8$ the kitchen is abandoned while the
bedroom is solved), so difficulty is geometry- and approach-dependent, not distance-dependent,
and no single flicker rate places all rooms in a partial band. This directed the aliasing rung
and is reported as calibration, not as a result.

\begin{figure}[h]
\centering
\includegraphics[width=0.52\linewidth]{p2_rung3_blackout_vs_p.png}
\caption{$P(\text{disabling blackout} \geq 1$ per episode$)$ versus flicker probability $p$ at
$k=4$, one curve per room using its real BFS path length $H$, with the hard-but-learnable band
shaded.}
\label{fig:blackout}
\end{figure}

\begin{figure}[h]
\centering
\includegraphics[width=0.52\linewidth]{p2_rung3_observed_vs_predicted.png}
\caption{Model-predicted (distance-$H$-based) per-room difficulty versus observed per-room
difficulty at $p=0.5/0.7/0.8$. The observed ordering does not follow $H$: difficulty is
geometry-, not distance-, driven.}
\label{fig:obspred}
\end{figure}

\paragraph{The ceiling reference.} $R(11)=27.8$ is the optimal return for the two near rooms,
which the trimmed IQM tracks; the full 16-D symmetric configuration itself sits just under it
at IQM 27.66. The ${\sim}1.6$ gap below that reference at $p=0.7$ is a path-length tax from
occasional blackouts, not sub-ceiling difficulty. At $p=0.7$, 9/10 seeds reach full success;
the tenth (seed 9) plateaus partway with bathroom SR 0.33, which the IQM trims but which
exists, and which is why that cell's per-room bathroom SR is 0.93 rather than 1.00.

\begin{figure}[h]
\centering
\includegraphics[width=0.70\linewidth]{p2_ladder_iqm.png}
\caption{The ladder cells of Table~\ref{tab:ladder} plotted: per-cell IQM eval return $+$ 95\%
CI. Four ceilings cluster near the 27.8 optimal-path reference; the $p=0.8$ fracture sits at
8.95 (200k) and its matched-budget fragmentation at 13.24 (500k), whose interval overlaps the
fracture's. The overlap is the point: the two budgets are indistinguishable in mean return and
differ only in per-seed distribution.}
\label{fig:ladder}
\end{figure}

\begin{figure}[h]
\centering
\includegraphics[width=0.92\linewidth]{p2_flicker08_perseed.png}
\caption{The $p=0.8$ fracture at the 200k budget, per seed. Left: outcome split (success /
timeout / collision), showing failures are timeouts, not collisions. Right: per-room
success-rate heatmap across the 10 seeds, kitchen abandoned 10/10 and bedroom solved 10/10.}
\label{fig:fracture}
\end{figure}

\paragraph{Per-seed outcomes at matched budget.} The fragmentation reported in the body is a
statement about the \emph{distribution} across seeds, so Table~\ref{tab:perseed} gives the
per-seed data it rests on. Two features stand out against the 200k result. First, the kitchen,
abandoned by every seed at 200k, is solved by five of ten seeds here, while the bedroom,
solved by every seed at 200k, is abandoned outright by seeds 2 and 6: no room is categorically
out of reach, and no room is reliably reached. Second, seeds 2 and 6 carry very low returns
(IQM 0.22 and 0.24) despite solving two of three rooms, because the sampled eval draw gave
them a majority of episodes on the room they abandoned. This illustrates why per-room SR and
aggregate return are reported separately rather than one being inferred from the other.

\begin{table}[h]
\caption{Per-seed outcomes for the $p=0.8$ flicker cell at the matched 500k budget.}
\label{tab:perseed}
\centering
\small
\begin{tabular}{llllll}
\toprule
Seed & kitchen SR & bedroom SR & bathroom SR & eval return IQM & regime \\
\midrule
0 & 0.00 & 1.00 & 0.33 & 15.70 & bedroom $+$ partial bathroom \\
1 & 0.00 & 1.00 & 0.67 & 17.96 & bedroom $+$ partial bathroom \\
2 & 1.00 & 0.00 & 1.00 & \phantom{0}0.22 & kitchen $+$ bathroom, bedroom abandoned \\
3 & 0.00 & 1.00 & 1.00 & 20.64 & bedroom $+$ bathroom, kitchen abandoned \\
4 & 1.00 & 1.00 & 1.00 & 25.56 & \textbf{all three rooms} \\
5 & 0.00 & 1.00 & 1.00 & 21.73 & bedroom $+$ bathroom, kitchen abandoned \\
6 & 1.00 & 0.00 & 1.00 & \phantom{0}0.24 & kitchen $+$ bathroom, bedroom abandoned \\
7 & 1.00 & 0.50 & 0.67 & 18.51 & kitchen $+$ partial bedroom/bathroom \\
8 & 0.00 & 1.00 & 0.67 & 17.76 & bedroom $+$ partial bathroom \\
9 & 1.00 & 1.00 & 1.00 & 25.13 & \textbf{all three rooms} \\
\midrule
\textbf{seed-mean} & \textbf{0.50} & \textbf{0.75} & \textbf{0.83} & \textbf{13.24} & fragmentation \\
\bottomrule
\end{tabular}
\end{table}

\clearpage
\section{Positive control in full}
\label{app:poscontrol}

The ladder applies a single fixed permutation of the MNIST pixels to a fraction
$q \in \{0, 0.25, 0.5, 0.75, 1.0\}$ of positions: nested subsets, a derangement within each,
identity elsewhere. Every rung is a bijection on pixel positions, so information is provably
preserved and the classes remain perfectly separable; what is progressively destroyed as
$q \to 1$ is only spatial locality. The CNN and the MLP are parameter-matched at 105.9k
parameters each, run at 10 seeds and a fixed sub-asymptotic budget where the convolutional
advantage is widest.

\begin{table}[h]
\caption{Positive-control dose-response: CNN versus parameter-matched MLP test accuracy
(IQM $+$ 95\% CI, 10 seeds) at each scramble fraction $q$.}
\label{tab:pc}
\centering
\small
\begin{tabular}{llll}
\toprule
$q$ & CNN test acc.\ (IQM, 95\% CI) & MLP test acc.\ (IQM, 95\% CI) & gap \\
\midrule
0.00 & 97.98 [97.81, 98.14] & 95.83 [95.72, 95.89] & \textbf{$+2.15$} \\
0.25 & 97.31 [97.12, 97.45] & 95.66 [95.54, 95.82] & $+1.65$ \\
0.50 & 96.27 [96.13, 96.47] & 95.69 [95.57, 95.86] & $+0.59$ \\
0.75 & 93.59 [93.25, 93.89] & 95.77 [95.68, 95.85] & $-2.18$ \\
1.00 & 93.18 [92.93, 93.45] & 95.70 [95.60, 95.82] & $-2.52$ \\
\bottomrule
\end{tabular}
\end{table}

\begin{figure}[h]
\centering
\includegraphics[width=0.55\linewidth]{pc_dose_response.png}
\caption{The dose-response of Table~\ref{tab:pc} plotted: CNN and MLP test accuracy versus
scramble fraction $q$, with 95\% CI bands. The CNN starts high, declines monotonically, and
crosses the flat MLP near $q\approx0.55$.}
\label{fig:pc}
\end{figure}

\paragraph{The frozen predictions.} \textbf{P1}, the CNN degrades monotonically as
$q \to 1$; \textbf{P2}, the MLP's achievable accuracy is identical, since the learning problem
is unchanged up to a relabeling of input units and the MLP hypothesis class is closed under
that relabeling; \textbf{P3}, the gap is large at $q=0$ and closes to ${\approx}0$ at $q=1$;
\textbf{P4}, the protocol therefore reports the task \emph{can} test $A$-suits-$S$ at $q=0$
and \emph{cannot} at $q=1$. P1 held (CNN IQM $97.98 \to 93.18$, monotone). P4 held. P3 was
partially wrong, in the informative direction reported in the body.

\paragraph{Why the MLP curve licenses the attribution.} Because each rung is a bijection, the
MLP's 0.13-point drift from $q=0$ to $q=1$ is optimisation noise, not information loss: the
task did not get harder, only locality was removed, so the CNN's fall can only be the loss of
the structure convolution exploits. Flatness is judged by a pre-committed absolute threshold,
$|\Delta| < 1.0$ point, which holds comfortably at 0.13. The corroborating check asks whether
the $q=1$ MLP IQM (95.70) lies inside the $q=0$ MLP's 95\% CI of [95.72, 95.89]; it misses at
the lower bound by 0.02 points. This is an artifact of precision, not a real shift: with 10
seeds that interval is only ${\pm}0.1$ point wide, so a change far smaller than any practical
significance still falls outside it. We pre-committed the absolute threshold precisely because
CI-overlap is not a sound flatness criterion at that resolution, which is pre-registration
working as intended. This is the step-5 attribution guard \emph{succeeding}, the complement of
the navigation $p=0.8$ rung where the same guard \emph{disqualified} a cell that had become
hard the wrong way.

\clearpage
\section{Why the task resists every manipulation}
\label{app:mechanism}

Three properties of the environment, each visible in the results, together explain why no
observation manipulation induced a hard-to-represent policy.

\paragraph{Observation redundancy.} The 16-D observation carries several channels that each
independently suffice for competent navigation, so removing any one leaves the others to carry
the load. The cumulative case is clearest: six of the sixteen dimensions can be removed,
global position (dims 9--10), distance-to-target (dim 11), and the region one-hot (dims
13--15), and the task is still solved completely, with all three rooms reached in 100\% of
episodes (IQM 27.17 [27.04, 27.37] at 10-D, against 27.66 [27.56, 27.78] for the full 16-D
observation). The half-point difference is statistically resolvable but is a path-length tax
rather than a loss of capability: the agent still reaches every target, taking marginally
longer. Nor is any single channel necessary in isolation. Removing global position alone
leaves a ceiling; additionally removing distance-to-target leaves a ceiling (27.24 [27.21,
27.27]); additionally removing the region one-hot leaves a ceiling at no detectable cost in
sample efficiency either (36.7k versus 40.0k environment steps to 90\% of asymptote, a
difference within seed noise). The proximity sensors appear to be the operative channel, since
they alone distinguish room, hallway, and doorway cells, which is what the region one-hot
nominally encoded. Yet even proximity's reliability is dispensable: corrupting it with 30\%
per-bit noise still ceilings (26.78 [26.41, 26.93]) at a collision rate of 0.00. The
observation is over-determined for this task.

\paragraph{Reactive optimality.} The optimal policy appears to be well approximated by a
function of the \emph{current} observation rather than of the observation history. This is why
frame-stacking neutralises flickering so effectively: with a stack of four frames, the agent
is deprived of every recent frame only when four consecutive maskings occur, an event of
probability $p^k$, which is 6.25\% at $p=0.5$ and 24\% at $p=0.7$. Because a single recent
frame suffices to act well, the agent need not integrate across gaps; it needs only
\emph{some} frame to be recent. Flickering therefore attacks the \emph{availability} of the
observation, not the \emph{complexity} of the function mapping observations to actions, and
availability is not the quantity the policy-hard hypothesis concerns. The distinction is
sharpest at $p=0.8$ (41\% fully masked), where behaviour no longer cleanly converges. At 200k,
returns fall to IQM 8.95 [6.43, 13.01] with a single dominant failure mode, kitchen
abandonment across every seed, failures dominated by timeouts (a per-seed timeout fraction of
0.13--0.47 of episodes) rather than collisions (${\approx}0$). At matched budget the aggregate
is not meaningfully higher (13.24 [7.79, 18.91]) but the per-seed structure fragments into six
distinct room-specialisation patterns, with two seeds finding a near-ceiling policy. Neither
budget shows the graceful degradation of a policy becoming harder to represent; both show an
optimisation landscape with multiple reactive local optima whose selection is seed-driven.
That is what the step-5 attribution guard exists to disqualify: hardness arising from
landscape structure over reactive policies, not from any policy being genuinely hard to
represent.

\paragraph{A single memorisable layout.} The map is fixed across all episodes, with optimal
path lengths of 11, 11, and 22 steps to the three targets. A policy can therefore encode
routes specific to this layout rather than computing them from the observation, which would
explain the robustness to proximity corruption, since a route-following policy needs little
reliable wall-sensing. We note this as an interpretation consistent with the data rather than
a directly measured property; we did not attempt to separate route memorisation from
observation-driven navigation, and doing so would require held-out layouts of the kind the
body's task specification calls for.

\paragraph{Why this closes off the observation axis.} The three properties compound. The
precondition the hypothesis requires, a policy hard to represent while the value function
remains smooth, is a property of the task's \emph{decision structure}, not of its observation
vector. This task's decision structure is short-horizon, static, and reactively solvable: the
optimal action is a simple function of local information, and it remains so however that
information is delivered. Degrading the observation can make the information scarcer
(flickering), narrower (removal), or noisier (corruption), but on this task none of these made
the underlying decision function more complex. Its complexity is set by the task's
short-horizon, static, reactively-solvable structure, not by how the observation is delivered.
Hence the pattern across the ladder: manipulations either leave the simple policy intact and
reachable (ceiling), or starve the agent enough that it stops attempting the task (fracture)
or scatters across seed-selected reactive optima (fragmentation). There is no intermediate
regime in which the policy becomes hard but learnable, because there is no hard policy to
find.

\clearpage
\section{Reproducibility}
\label{app:repro}

Every quantitative claim in this paper is backed by committed per-seed data and regenerable
analysis. Every aggregate reported, including every confidence interval, regenerates from the
released per-seed CSVs under a seeded bootstrap: the resampling seed is fixed in the analysis
scripts, so intervals reproduce to the digit rather than wobbling between runs. We verified
this by re-deriving every claim-grade number in the paper from the released CSVs, covering all
six ladder cells, the Phase 1 aggregates and probability of improvement, steps-to-90\% for all
nine configurations, the redundancy and timeout figures, every per-seed row, and all ten
positive-control cells.

The supplementary material contains: (i) the per-seed raw CSVs behind every IQM and confidence
interval; (ii) the analysis and plotting scripts that turn those CSVs into every figure and
table, using rliable for IQM aggregation with stratified-bootstrap confidence intervals;
(iii) the frozen pre-registration documents for both the observability ladder and the positive
control, including the dated amendment trail; (iv) the environment and observation-wrapper
code implementing the frozen navigation task and the four degradation mechanisms; and (v) a
running results log recording each experiment's headline numbers and interpretation in the
order they were obtained.

The pre-registration documents fix, before any run, the hypotheses, the ladder, the
falsification criteria, and the per-outcome interpretation table on which the methodological
claim of this paper rests. They are included so that the ordering of design decisions relative
to results can be independently checked.

\clearpage
\section{NeurIPS paper checklist}
\label{app:checklist}

%%%% NeurIPS paper checklist. Answers reflect the submitted version only.
%%%% Anonymity: no answer may name an author, institution, or repository.

\begin{enumerate}

\item {\bf Claims}\\
Do the main claims made in the abstract and introduction accurately reflect the paper's
contributions and scope?
\answerYes{}
\begin{itemize}
\item[] Justification: The abstract and introduction claim exactly two things, a reusable
  pre-registered protocol and a worked negative outcome, and both are delivered in
  \S\ref{sec:protocol} and \S\ref{sec:results}. The abstract states explicitly that the
  result is a statement about the task rather than evidence against the architecture
  hypothesis, which is the paper's central scope limit.
\end{itemize}

\item {\bf Limitations}\\
Does the paper discuss the limitations of the work performed by the authors?
\answerYes{}
\begin{itemize}
\item[] Justification: \S\ref{sec:discussion} states the scope limit directly: we show this
  task class cannot test the hypothesis, and explicitly have \emph{not} shown that inverted
  asymmetry fails in reinforcement learning generally. \S\ref{sec:poscontrol} reports a
  pre-registered prediction (P3) that was wrong, and Appendix~\ref{app:mechanism} flags the
  memorisation account as an interpretation consistent with the data rather than a directly
  measured property. The positive control is on image classification, not RL, which is stated
  as a limitation rather than glossed.
\end{itemize}

\item {\bf Theory assumptions and proofs}\\
Does the paper provide the full set of assumptions and a complete (and correct) proof?
\answerNA{}
\begin{itemize}
\item[] Justification: The paper is empirical and methodological; it contains no theoretical
  results or proofs.
\end{itemize}

\item {\bf Experimental result reproducibility}\\
Does the paper fully disclose all the information needed to reproduce the main experimental
results?
\answerYes{}
\begin{itemize}
\item[] Justification: Appendix~\ref{app:repro} describes the full pipeline. All
  hyperparameters live in configuration files included in the supplementary material, the
  environment and observation-wrapper code is included, training is seeded end-to-end, and the
  bootstrap resampling seed is fixed so that every reported confidence interval reproduces to
  the digit from the released per-seed CSVs.
\end{itemize}

\item {\bf Open access to data and code}\\
Does the paper provide open access to the data and code, with sufficient instructions to
faithfully reproduce the main experimental results?
\answerYes{}
\begin{itemize}
\item[] Justification: An anonymised supplementary archive accompanies the submission
  containing the per-seed CSVs, the analysis and plotting scripts, the environment and wrapper
  code, the configuration files, and the frozen pre-registration documents, with reproduction
  commands in its README. A public repository link will be added in the camera-ready version.
\end{itemize}

\item {\bf Experimental setting/details}\\
Does the paper specify all the training and test details necessary to understand the results?
\answerYes{}
\begin{itemize}
\item[] Justification: \S\ref{sec:phase1} gives the environment, action space, observation
  dimension, reward function, and step budget. Appendix~\ref{app:phase1} gives the full
  configuration table with per-configuration network architectures and the varied
  hyperparameters. Appendix~\ref{app:phase2} gives the ladder mechanisms and their parameters;
  every remaining hyperparameter is in the supplementary configuration files.
\end{itemize}

\item {\bf Experiment statistical significance}\\
Does the paper report error bars suitably and correctly defined or other appropriate
information about the statistical significance of the experiments?
\answerYes{}
\begin{itemize}
\item[] Justification: Every aggregate is an interquartile mean with a 95\% stratified
  bootstrap confidence interval over 10 seeds, computed with rliable. The aggregation
  conventions, including the seed-mean definition for per-room success rates and the exclusion
  of a sentinel value from the sample-efficiency statistic, are stated in
  Appendix~\ref{app:phase1}. Single-seed results are used only as learnability gates and are
  never reported as findings, a distinction the paper states explicitly and enforces.
\end{itemize}

\item {\bf Experiments compute resources}\\
Does the paper provide sufficient information on the computer resources needed to reproduce
the experiments?
\answerYes{}
\begin{itemize}
\item[] Justification: All runs were executed serially on a single commodity cloud GPU
  instance. Each Phase 1 configuration is 10 seeds at 200k environment steps; the Phase 2
  ladder adds six 10-seed cells at 200k or 500k steps. The positive control is CPU-feasible at
  10 seeds and a two-epoch budget. Per-run step budgets are given in
  Table~\ref{tab:ladder} and Appendix~\ref{app:phase1}.
\end{itemize}

\item {\bf Code of ethics}\\
Does the research conducted in the paper conform, in every respect, with the NeurIPS Code of
Ethics?
\answerYes{}
\begin{itemize}
\item[] Justification: The work uses a synthetic grid-world environment and MNIST. It involves
  no human subjects, no personal data, and no deployed system.
\end{itemize}

\item {\bf Broader impacts}\\
Does the paper discuss both potential positive societal impacts and negative societal impacts
of the work performed?
\answerNA{}
\begin{itemize}
\item[] Justification: The contribution is a methodological check on whether a benchmark can
  test a hypothesis. Its foreseeable effect is on evaluation practice within machine learning
  research, and we identify no direct societal impact, positive or negative, beyond improved
  reliability of published empirical claims.
\end{itemize}

\item {\bf Safeguards}\\
Does the paper describe safeguards that have been put in place for responsible release of data
or models that have a high risk for misuse?
\answerNA{}
\begin{itemize}
\item[] Justification: The released artifacts are small grid-world policies, analysis scripts,
  and per-seed result tables. They carry no misuse risk.
\end{itemize}

\item {\bf Licenses for existing assets}\\
Are the creators or original owners of assets used in the paper properly credited and are the
license and terms of use explicitly mentioned and properly respected?
\answerYes{}
\begin{itemize}
\item[] Justification: The paper cites the rliable evaluation methodology, the PPO algorithm,
  and MNIST-based experimental design, and the supplementary material lists all software
  dependencies with their versions. All are standard open-source research assets used in
  accordance with their licenses.
\end{itemize}

\item {\bf New assets}\\
Are new assets introduced in the paper well documented and is the documentation provided
alongside the assets?
\answerYes{}
\begin{itemize}
\item[] Justification: The supplementary archive documents the environment, the four
  observation-degradation wrappers, the analysis scripts, and the per-seed data, with a README
  giving the layout, a data-to-claim map linking each paper element to the CSV and script that
  produce it, and reproduction commands.
\end{itemize}

\item {\bf Crowdsourcing and research with human subjects}\\
Does the paper include the full text of instructions given to participants and screenshots, if
applicable, as well as details about compensation?
\answerNA{}
\begin{itemize}
\item[] Justification: The paper involves no crowdsourcing and no human subjects.
\end{itemize}

\item {\bf Institutional review board (IRB) approvals}\\
Does the paper describe potential risks incurred by study participants, disclose whether such
risks were disclosed to the subjects, and obtain IRB approvals?
\answerNA{}
\begin{itemize}
\item[] Justification: The paper involves no human subjects, so no IRB review was required.
\end{itemize}

\item {\bf Declaration of LLM usage}\\
Does the paper describe the usage of LLMs if it is an important, original, or non-standard
component of the core methods in this research?
\answerNA{}
\begin{itemize}
\item[] Justification: No large language model forms any part of the core method. The
  contribution is an experimental protocol and its instantiation on reinforcement-learning and
  image-classification tasks.
\end{itemize}

\end{enumerate}

